\documentclass[11pt]{article}
\usepackage[T1]{fontenc}
\usepackage[utf8]{inputenc}
\usepackage{lmodern}
\usepackage[a4paper,margin=0.88in]{geometry}
\usepackage{amsmath,amssymb,amsthm,mathtools}
\usepackage{microtype}
\usepackage{booktabs,tabularx,longtable,array}
\usepackage{enumitem}
\usepackage{xcolor}
\usepackage[hidelinks]{hyperref}
\usepackage[nameinlink,capitalize]{cleveref}
\usepackage{fancyhdr}
\usepackage{lastpage}

\IfFileExists{buildmeta.tex}{\input{buildmeta.tex}}{%
  \newcommand{\BuildCommit}{local-audit}%
  \newcommand{\BuildDate}{28 August 2026}%
}

\definecolor{deepblue}{RGB}{25,55,92}
\definecolor{softblue}{RGB}{235,242,249}
\definecolor{softgray}{RGB}{245,245,245}
\definecolor{softgreen}{RGB}{237,247,239}
\definecolor{softred}{RGB}{252,239,239}
\definecolor{softamber}{RGB}{255,247,230}

\newtheorem{theorem}{Theorem}[section]
\newtheorem{proposition}[theorem]{Proposition}
\newtheorem{lemma}[theorem]{Lemma}
\newtheorem{corollary}[theorem]{Corollary}
\theoremstyle{definition}
\newtheorem{definition}[theorem]{Definition}
\newtheorem{remark}[theorem]{Remark}

\newcommand{\N}{\mathbb N}
\newcommand{\Q}{\mathbb Q}
\newcommand{\R}{\mathbb R}
\newcommand{\C}{\mathbb C}
\newcommand{\1}{\mathbf 1}
\newcommand{\E}{\mathcal E}
\newcommand{\B}{\mathcal B}
\newcommand{\WPS}{\operatorname{WindowPairSupply}}
\newcommand{\status}[1]{\textsf{#1}}
\newcommand{\lean}{\textsc{Lean}}

\pagestyle{fancy}
\fancyhf{}
\lhead{Erd\H{o}s Problem \#287 -- Clean Public Review V16}
\rhead{28 August 2026}
\cfoot{\thepage\ of \pageref{LastPage}}
\setlength{\headheight}{14pt}
\setlength{\emergencystretch}{2em}

\title{\vspace{-1.2cm}\textbf{A Machine-Checked Finite Reduction for Erd\H{o}s Problem \#287\\
and a Source-Explicit Large-$M$ Criterion}\\[1ex]
\large Clean Public Review -- Version 16}
\author{Research draft with Lean/Aristotle-assisted formal verification}
\date{28 August 2026}

\begin{document}
\maketitle

\begin{center}
\fcolorbox{deepblue}{softblue}{\begin{minipage}{0.92\textwidth}
\textbf{Status.} Erd\H{o}s Problem \#287 remains open.  Public repository \texttt{main}
machine-checks the exact problem statement, the finite exclusion through maximum denominator
$4\times10^9$, the conditional window-pair closure compiler, and the V14 finite/algebraic
Vaughan source spine.  The later V15 $\Lambda=\mu*\log$ and balanced-seven Lean files are
present in the supplied midpoint Aristotle archive, but were not yet on public \texttt{main}
at this audit cutoff and were not independently rebuilt in the present execution environment.
They are therefore labelled \emph{formalization source present -- public replay pending}, not
public Lean theorems.  No prime-distribution estimate, Ford--Maynard lower bound, or final
large-$M$ supply is claimed proved.  V16 source commit: \texttt{\BuildCommit}.
\end{minipage}}
\end{center}

\begin{abstract}
Erd\H{o}s Problem \#287 asks whether every representation
\[
1=\frac1{n_1}+\cdots+\frac1{n_k},\qquad 1<n_1<\cdots<n_k,
\]
contains a consecutive gap of size at least $3$.  This note separates the present work into
three logically distinct layers.

First, the finite layer is machine-checked: there is no counterexample whose maximum
denominator is at most $4\times10^9$.  The proof uses a prime-power window obstruction,
an adjacent-hole contradiction, and a chain of interval certificates rather than a search
over all unit-fraction representations.

Second, the large-$M$ conclusion is compiled from one explicit window-pair supply predicate.  It asks for two consecutive integers in the top half of $[1,M]$, each carrying
a sufficiently large prime-power divisor.  A Lean theorem shows that an effective threshold
$M_0\le4\times10^9$ together with $\WPS(M)$ for every $M\ge M_0$ implies the exact public
statement.  No inhabitant of this large-$M$ predicate is presently known.

Third, the current analytic research is made source-explicit.  The public V14 bank
contains an exact Vaughan source spine.  A supplied midpoint V15 formalization develops the
source-minimal identity $\Lambda=\mu*\log$, its affine specialization at $2mn\pm1$, an
exact three-way factor partition, a determinant-one line parametrization, and finite
balanced-seven polarization algebra.  Those V15 files are treated here as formalization
source pending final public replay; the mathematical identities are displayed explicitly
and do not depend on an analytic estimate.  The remaining analytic statements are a signed
endpoint distribution estimate for a polarized family of $X^{1/7}$-smooth multiplicative
functions at moduli $X^{3/5}$, its comparison/low-conductor matching statement, and the
remaining source-minimal hard-line estimates.  These are open specifications, not theorems.
\end{abstract}

\section*{Status at a glance}
\addcontentsline{toc}{section}{Status at a glance}
\begin{center}
\renewcommand{\arraystretch}{1.22}
\begin{tabularx}{\textwidth}{>{\raggedright\arraybackslash}p{0.31\textwidth}
>{\raggedright\arraybackslash}p{0.19\textwidth}X}
\toprule
Item & Status & Meaning \\
\midrule
Exact public counterexample predicate & \status{PUBLIC LEAN-CHECKED} & Set form and ordered sequence form are connected without adding a hidden $k\ge2$ hypothesis. \\
Finite range $M\le4\times10^9$ & \status{PUBLIC LEAN-CHECKED} & A 34-certificate interval chain excludes every exact counterexample in this range. \\
Window-pair closure compiler & \status{PUBLIC LEAN-CHECKED, CONDITIONAL} & An explicit large-$M$ supply predicate implies the full conjecture when paired with the finite range. \\
V14 Vaughan affine source spine & \status{PUBLIC LEAN-CHECKED ALGEBRA} & Exact decomposition and structural compilers; analytic inputs remain uninhabited. \\
$\Lambda=\mu*\log$, hard-pair line, balanced-seven algebra & \status{FORMALIZATION REPLAY PENDING} & V15 source is present in the midpoint archive and statically audited; it was not yet replayed on public \texttt{main} at cutoff. \\
Large-$M$ prime distribution & \status{OPEN} & No proof of $\WPS(M)$ for every sufficiently large $M$. \\
Polarized signed endpoint estimate & \status{OPEN ANALYTIC} & Current exact V-branch residual at $(Q,y)=(X^{3/5},X^{1/7})$. \\
Comparison/low-conductor matching & \status{OPEN SOURCE} & The physical comparison sequence must undergo the same decomposition. \\
Erd\H{o}s Problem \#287 & \status{OPEN} & No unconditional proof is claimed. \\
\bottomrule
\end{tabularx}
\end{center}

\paragraph{Status vocabulary.}
\status{PUBLIC LEAN-CHECKED} means the declaration is on public repository
\texttt{main}.  \status{EXACT ALGEBRA} means the identity is proved in the exposition.
\status{FORMALIZATION REPLAY PENDING} means Lean source is present in the supplied midpoint
archive but has not yet been promoted into the public bank.  \status{CONDITIONAL} means a
compiler is proved from an explicit antecedent.  \status{OPEN} means that antecedent or
analytic estimate is not proved.  \status{FALSE/RETIRED} is reserved for a literal failed
dictionary or superseded route, never for the underlying general theorem.

\tableofcontents

\section{The problem and the scope of this note}

The public problem asks whether, for distinct integers
\[
1<n_1<\cdots<n_k,\qquad \sum_{i=1}^k\frac1{n_i}=1,
\]
one necessarily has
\[
\max_{1\le i<k}(n_{i+1}-n_i)\ge3.
\]
It is listed as open in the Erd\H{o}s Problems archive and originates in the sources
cited there \cite{ErdosGraham1980,ErdosProblems287}.

For a finite set $A\subset\N$, write
\[
N=\min A,\qquad M=\max A.
\]
A counterexample is a finite set satisfying
\[
A\subset\{2,3,\dots\},\qquad \sum_{a\in A}\frac1a=1,
\]
and the local gap condition that every nonmaximal $a\in A$ is followed by $a+1$ or $a+2$
in $A$.  The Lean development formalizes exactly this predicate and proves the ordered
form from it.

This note is deliberately not a chronological research diary.  It records only the
current controlling finite theorem, the exact conditional closure interface, the
source algebra with its precise public/formalization status, and the current open analytic
obligations.  Earlier
provider searches, hostile audits, and superseded residuals are retained in the separate
technical verification ledger and immutable V15 archive.

\section{The machine-checked finite theorem}

\begin{theorem}[Finite exclusion]
There is no exact counterexample to Erd\H{o}s Problem \#287 with
\[
3\le M\le4{,}000{,}000{,}000.
\]
Equivalently, every unit-fraction representation in the problem with maximum denominator
at most $4\times10^9$ has a consecutive gap at least $3$.
\end{theorem}

The theorem is represented in the repository by
\texttt{no\_Erdos287Counterexample\_of\_max\_le\_4e9}.  Its proof is not an exhaustive
enumeration of finite subsets.  It is a certificate compiler built from the three ideas
below.

\subsection{Holes and the adjacent-hole contradiction}

For a hypothetical counterexample, call an integer $x\in[N,M]$ a \emph{hole} when
$x\notin A$.  The gap-$\le2$ condition implies that two consecutive integers in
$[N,M]$ cannot both be holes.  This is the formal theorem
\texttt{Gap2CE.holes\_isolated}.

A second elementary fact places the top half of the interval above the minimum:
\[
N\le\left\lfloor\frac M2\right\rfloor.
\]
Thus, whenever $M\le2x$ and $x+1\le M$, the consecutive pair $x,x+1$ lies inside the
closed window where the adjacent-hole contradiction applies.

\subsection{Prime-power windows}

For $j\ge0$, let $C(j)$ denote the largest reduced numerator arising from a nonempty
partial reciprocal sum
\[
\sum_{s\in S}\frac1s,\qquad \varnothing\ne S\subseteq\{1,\dots,j\}.
\]
The formal development uses a certified upper table
\[
C_{\mathrm{val}}(0),\dots,C_{\mathrm{val}}(9)
=1,1,3,11,25,137,137,1019,2143,7129.
\]

The prime-power window theorem says, schematically, that if $p$ is prime, $e\ge1$,
$p^e\mid x$, and
\[
C\!\left(\left\lfloor\frac M{p^e}\right\rfloor\right)<p,
\]
then $x$ cannot belong to $A$.  The reason is that the top $p$-adic layer of the reciprocal
sum would force a nonzero numerator to vanish modulo $p$.

Consequently, if both $x$ and $x+1$ possess such prime-power divisors, then they are two
adjacent holes.  This gives the central finite compiler.

\begin{proposition}[Window-pair blocker]
Suppose there are primes $p_u,p_v$, positive exponents $a_u,a_v$, and an integer $x$ such
that
\[
p_u^{a_u}\mid x,\qquad p_v^{a_v}\mid x+1,
\]
\[
\left\lfloor\frac M{p_u^{a_u}}\right\rfloor,
\left\lfloor\frac M{p_v^{a_v}}\right\rfloor\le9,
\]
\[
C_{\mathrm{val}}\!\left(\left\lfloor\frac M{p_u^{a_u}}\right\rfloor\right)<p_u,
\qquad
C_{\mathrm{val}}\!\left(\left\lfloor\frac M{p_v^{a_v}}\right\rfloor\right)<p_v,
\]
and
\[
M\le2x,
\qquad x+1\le M.
\]
Then a gap-$\le2$ counterexample with maximum $M$ cannot exist.
\end{proposition}

\subsection{The interval-certificate chain}

One certificate can exclude an entire interval $L\le M\le U$.  The divisibility data are
fixed at a chosen $x$, monotonicity controls the windows as $M$ varies, and $U\le2x$
keeps the forced holes above the minimum.  The repository contains a chain of $34$
certificates whose intervals meet end-to-end and cover
\[
[3,4\times10^9].
\]
All primality side conditions are discharged in the kernel; the construction does not use
\texttt{native\_decide}.

\section{The exact large-$M$ closure criterion}

\begin{definition}[Window-pair supply]
For $M\in\N$, $\WPS(M)$ is the assertion that there exist
$x,p_u,a_u,p_v,a_v\in\N$ such that
\[
\begin{gathered}
p_u,p_v\text{ are prime},\qquad a_u,a_v\ge1,\\
p_u^{a_u}\mid x,\qquad p_v^{a_v}\mid x+1,\\
\left\lfloor\frac M{p_u^{a_u}}\right\rfloor\le9,
\quad C_{\mathrm{val}}\!\left(\left\lfloor\frac M{p_u^{a_u}}\right\rfloor\right)<p_u,\\
\left\lfloor\frac M{p_v^{a_v}}\right\rfloor\le9,
\quad C_{\mathrm{val}}\!\left(\left\lfloor\frac M{p_v^{a_v}}\right\rfloor\right)<p_v,\\
M\le2x,
\qquad x+1\le M.
\end{gathered}
\]
\end{definition}

In plain terms, two consecutive integers in the top half of $[1,M]$ each have a
prime-power divisor larger than roughly $M/10$, with a small local numerator condition.

\begin{theorem}[End-to-end conditional compiler]
Assume there is an explicit $M_0\le4\times10^9$ such that
\[
\WPS(M)\qquad\text{for every }M\ge M_0.
\]
Then Erd\H{o}s Problem \#287 is true.
\end{theorem}

This implication is machine-checked as
\texttt{no\_Erdos287Counterexample\_of\_closure}.  The proof splits at
$4\times10^9$: the finite theorem handles the lower range, while the window-pair blocker
handles every larger maximum.

The interface is weaker than a Sophie-Germain prime condition.  For example, if a prime
$q$ in the interval $M/3<q\le(M-1)/2$ has $2q+1$ prime, then one may take
$x=2q$.  The analogous minus case uses $x=2q-1$.  The Lean theorem
\texttt{windowPairSupply\_of\_sophieWitness} formalizes this implication.  The present
programme does not claim such Sophie witnesses for every large $M$.

\section{The source-minimal affine algebra}

This section concerns the analytic route toward a large-$M$ supply.  The formulas below
are exact and are proved algebraically in this note.  Matching V15 Lean source is present in
the supplied midpoint archive, but final public replay was pending at the audit cutoff;
accordingly this section is not labelled public Lean-checked.  All cancellation estimates
remain open.

\subsection{The identity $\Lambda=\mu*\log$}

For the genuine arithmetic functions, Dirichlet convolution gives
\[
\mu*\log=\mu*(\Lambda*1)=(\mu*1)*\Lambda=\Lambda,
\]
and hence
\[
\boxed{\Lambda(N)=\sum_{qr=N}\mu(q)\log r.}
\]
Equivalently,
\[
\Lambda(N)=\sum_{q\mid N}\mu(q)\log\frac Nq.
\]
The proof is carried out in the Dirichlet convolution ring from
$\mu*1=\varepsilon$ and $\log=\Lambda*1$; it agrees definitionally with Mathlib's own
identity.

At the affine value
\[
N=2mn+s,
\qquad s\in\{-1,+1\},
\]
one obtains the exact source
\[
\Lambda(2mn+s)=\sum_{qr=2mn+s}\mu(q)\log r.
\]
The sign is represented by a dedicated formal datatype, preventing illegal subtraction in
$\N$ for the case $s=-1$.

\subsection{Exact factor partition}

For a cutoff $U$, the factor pairs are partitioned disjointly into
\[
q\le U,
\qquad q>U,\ r\le U,
\qquad q>U,\ r>U.
\]
The three classes are exhaustive and pairwise disjoint, so an arbitrary weighted sum
recombines exactly over them.  The hard part is therefore
\[
\sum_{\substack{qr=2mn+s\\q>U,\ r>U}}\mu(q)\log r.
\]
No estimate for this sum is encoded.

\subsection{The determinant-one line}

Every hard pair satisfies
\[
qr-2mn=s,
\qquad s=\pm1.
\]
This unit shift forces
\[
(r,2m)=1,
\qquad(q,2n)=1.
\]
Fixing $r,m$ and one solution $(q_0,n_0)$, all integer solutions are exactly
\[
\boxed{q=q_0+2mt,
\qquad n=n_0+rt,
\qquad t\in\mathbb Z,}
\]
and the parameter $t$ is unique.  The forward direction, reverse direction, coprimality,
nonvanishing, and multiplicity one are elementary consequences of the unit determinant;
matching V15 Lean declarations are present in the midpoint archive pending public replay.

The exact exponent ledger records, for the programme parameter
$\sigma<1/6$,
\[
\sigma+\frac23+\frac16<1,
\qquad
\frac16-\frac{16623}{100000}=\frac{131}{300000}>0.
\]
These are exponent identities only; they are not bounds for an analytic sum.

\begin{remark}[Relation to the V14 Vaughan route]
The earlier exact Vaughan decomposition and prime-modulus M\"obius two-outer packet remain
valid alternative source decompositions.  They are not false, but the identity
$\Lambda=\mu*\log$ is source-minimal and does not require truncation parameters or an
additional convolution with $1$.  The clean public paper therefore treats the Vaughan
packet as an optional analytic child, not as the unique controlling source.
\end{remark}

\section{The balanced-seven polarization}

One explicit V-branch cell has seven labelled prime boxes.  Put
\[
Q=X^{3/5},
\qquad Y=X^{1/7},
\]
and consider
\[
\B_{7,s}(X)=
\sum_{q\sim Q}\mu(q)
\left[
\sum_{\substack{p_1,\dots,p_7\sim Y\\2p_1\cdots p_7\equiv-s\pmod q}}
\prod_{i=1}^7\omega_i(p_i)
-\mathfrak M_{7,s}(q)
\right].
\]
The desired scale is $o(X/\log X)$.

\subsection{Why separate character energies do not close}

Character orthogonality produces a sevenfold character product.  Splitting it $3+4$ and
using the ordinary multiplicative large sieve gives the scale
\[
Q\,(X^{3/7}X^{4/7})^{1/2}
=X^{3/5+1/2}
=X^{11/10},
\]
a deficit of $X^{1/10}$.  Separate high-moment interpolation gives the same scale.  This
explains why the current reduction does not simply repeat a standard large-sieve argument.

\subsection{Exact squarefree polarization}

For $\mathbf z=(z_1,\dots,z_7)\in\mathbb T^7$, define
\[
a_{\mathbf z}(p)=\frac17\sum_{i=1}^7z_i\omega_i(p).
\]
Extend this to a squarefree multiplicative function $f_{\mathbf z}$ supported on primes in
the box $p\sim Y$ by
\[
f_{\mathbf z}(p)=a_{\mathbf z}(p),
\qquad f_{\mathbf z}(p^\ell)=0\quad(\ell\ge2).
\]
Expanding the polynomial shows that the coefficient of $z_1\cdots z_7$ in
\[
\prod_{j=1}^7\left(\sum_{i=1}^7z_i\omega_i(p_j)\right)
\]
is exactly
\[
\sum_{\sigma\in S_7}\prod_{j=1}^7\omega_{\sigma(j)}(p_j),
\]
with each permutation appearing once.  The normalization factor is
$7^{-7}=1/823543$, and $|S_7|=5040$.

Repeated-prime tuples have only six free prime variables, hence contribute at most
\[
X^{6/7+o(1)},
\]
which is below $X/\log^A X$ for every fixed $A$.

The same source audit gives the finite squarefree encoding and the binomial certificate
\[
\sum_{j=0}^{3}(-1)^j\binom7j=-20.
\]
It does not identify this number with a physical Ford weight; the statement is purely
finite combinatorics.

\subsection{The exact V-branch decomposition}

Let $\Delta_{\mathrm{LC}}(f_{\mathbf z};X,q,a_s)$ denote the progression discrepancy after
removing the chosen principal and low-conductor terms, where
$a_s\equiv-s\overline2\pmod q$.  Coefficient extraction gives
\[
\begin{aligned}
\B_{7,s}(X)
={}&7^7\int_{\mathbb T^7}\overline{z_1\cdots z_7}
\sum_{q\sim X^{3/5}}\mu(q)
\Delta_{\mathrm{LC}}(f_{\mathbf z};X,q,a_s)\,d\mathbf z\\
&+\sum_{q\sim Q}\mu(q)
\bigl(\mathfrak M^{\mathrm{pol}}_{7,s}(q)-\mathfrak M_{7,s}(q)\bigr)
+O(X^{6/7+o(1)}).
\end{aligned}
\]
Thus this cell is reduced to two source-exact statements.

\paragraph{Polarized signed endpoint estimate.}
\[
\int_{\mathbb T^7}\overline{z_1\cdots z_7}
\sum_{q\sim X^{3/5}}\mu(q)
\Delta_{\mathrm{LC}}(f_{\mathbf z};X,q,a_s)\,d\mathbf z
=o(X/\log X).
\]
Programme label:
\path{AFFINE287-POLARIZED-OMEGA7-SIGNED-EOD45}.

\paragraph{Comparison/low-conductor match.}
\[
\sum_{q\sim Q}\mu(q)
\bigl(\mathfrak M^{\mathrm{pol}}_{7,s}(q)-\mathfrak M_{7,s}(q)\bigr)
=o(X/\log X).
\]
Programme label:
\path{AFFINE287-MULOG-COMPARISON-LOWCOND-MATCH45}.

Both are open.  The supplied V15 status source defines uninhabited specification structures
for them, not proofs; that status file was not yet part of public \texttt{main} at the audit
cutoff.

\section{What existing distribution theorems do and do not supply}

Ford and Maynard provide a general prime-producing sieve framework from Type I and Type II
information for an actual nonnegative sequence \cite{FordMaynard2024}.  Their theorem is a
downstream framework, not a proof of the specific affine correlations above.

Drappeau, Granville and Shao prove distribution of smooth-supported multiplicative
functions beyond the square-root barrier, up to moduli $x^{3/5-\varepsilon}$ under their
stated hypotheses \cite{DGS2017}.  Pascadi obtains stronger exponents of distribution for
primes and smooth numbers in a different weighted setting \cite{Pascadi2025}.  These
results are close in scale to the polarized cell, but the current audit has not established
a literal specialization at the endpoint
\[
(Q,y)=(X^{3/5},X^{1/7})
\]
with the signed M\"obius modulus average, torus coefficient, low-conductor subtraction, and
physical comparison term all preserved.  This note therefore records a promising
neighbourhood of existing technology, not a black-box closure.

Several negative conclusions are also retained:
\begin{itemize}[leftmargin=1.6em]
\item the ordinary separate multiplicative large sieve misses by $X^{1/10}$;
\item a three-block decomposition satisfying the relevant asymmetric ranges is unavailable
for seven equal $1/7$ exponents;
\item a termwise prime decomposition cannot remove every unsplit prime-box term;
\item a modulus micro-switch cannot be exhaustive because prime moduli survive;
\item the V14 truncated M\"obius cofactor is not literally a well-factorable modulus weight;
\item the canonical Gate-1A determinant-$2k$ row does not directly match a determinant
$\pm1$ affine source.
\end{itemize}
The detailed proofs and route-by-route provenance belong in the technical ledger, not in
the main paper.

\section{The current end-to-end dependency graph}

The verified and open layers can now be displayed without historical notation:
\[
\boxed{\begin{array}{c}
\text{exact problem statement}\\
\downarrow\\
\text{finite window compiler and certificates}\\
\downarrow\\
M\le4\times10^9\text{ excluded}
\end{array}}
\qquad
\boxed{\begin{array}{c}
\text{large-}M\ \WPS(M)\text{ supply}\\
\downarrow\\
\text{public Lean closure compiler}\\
\downarrow\\
\text{Erd\H{o}s \#287}
\end{array}}
\]
The analytic programme seeks to produce the missing large-$M$ supply.  Its current
source-explicit front is
\[
\boxed{\begin{array}{c}
\Lambda=\mu*\log\text{ affine source}\\
+\ \text{balanced-seven polarization}\\
\downarrow\\
\text{hard affine-line cancellation}\\
+\ \text{polarized signed endpoint distribution}\\
+\ \text{comparison matching}\\
\downarrow\\
\text{positive affine prime/almost-prime mass}\\
\downarrow\\
\WPS(M)\text{ or a stronger Sophie-type witness}
\end{array}}
\]
Every downward arrow in the second display contains open analysis.  The diagram is a
research programme, not a proof.  Moreover, closing the displayed balanced-seven cell does
not by itself close all remaining fixed-certificate packets or the full large-$M$ supply.

\section{Formal verification and reproducibility}

\subsection{Public Lean bank}
At V16 source commit \texttt{\BuildCommit}, public repository \texttt{main} contains the
exact problem predicate, finite exclusion through $4\times10^9$, the window-pair closure
compiler, and the V14 Vaughan structural spine.  The public V14 status explicitly leaves
\path{AFFINE287-PRIME-MODULUS-MU-TWOOUTER45}, the $k=0$ smooth packet,
fixed-certificate leakage, Gate closure, and the large-$M$ supply open.  The accompanying
V14 report records a successful \texttt{lake build} of $8111$ jobs and no Lean-code use of
\texttt{sorry}, \texttt{admit}, a user \texttt{axiom}, \texttt{opaque},
\texttt{unsafe}, \texttt{native\_decide}, or \texttt{@[implemented\_by]}.

\subsection{Midpoint V15 archive}
The supplied \texttt{midway.gz} archive contains $92$ Lean files, including the
$\Lambda=\mu*\log$ and balanced-seven modules and a V15 status file.  A static audit found
no Lean-code occurrences of the prohibited declarations above, and the two analytic
interfaces are deliberately uninhabited.  However:
\begin{itemize}[leftmargin=1.6em]
\item those V15 files were not yet present on public \texttt{main} at the audit cutoff;
\item the archive's top-level build summary records the completed V14 replay, while the user
reported that the V15 Lean run was still in progress;
\item the present execution environment did not contain \texttt{lake}, so no independent
local V15 kernel replay was performed.
\end{itemize}
V16 therefore does not promote the V15 formalization to \status{PUBLIC LEAN-CHECKED}.
Once the ongoing replay completes and is merged, the formal-status table can be promoted
without changing the mathematical exposition.

\subsection{Document organization}
The previous V15 PDF is preserved as an immutable technical patch-stack archive.  This
V16 document is the coherent public paper; hostile audits, failed provider dictionaries,
historical residuals, Gate-specific archaeology, and full chronological provenance live in
the separate technical verification ledger.  The clean source is compiled reproducibly in
the public repository, and the legacy forum PDF path is retained as an alias to the current
clean paper.

\section{Conclusion}

The strongest unconditional result presently available in this project is finite:
Erd\H{o}s Problem \#287 has no counterexample with maximum denominator at most
$4\times10^9$.  The exact remaining global input is also clean: prove $\WPS(M)$ for every
sufficiently large $M$ with an effective threshold inside the verified finite range.

The recent analytic work has reduced parts of the source without claiming cancellation.
The identity $\Lambda=\mu*\log$, the hard affine determinant-one line, and the
balanced-seven polarization are now exact algebraic infrastructure with V15 formalization
source present and final public replay pending.  The next
mathematics is not another document-level reorganization: it is a genuine signed
prime-distribution estimate and a source-matched comparison theorem.  Until those are
proved and compiled into a large-$M$ supply, Erd\H{o}s Problem \#287 remains open.

\appendix
\section{Formal theorem and source map}

\subsection*{Public Lean bank at the V16 cutoff}
\begin{itemize}[leftmargin=1.6em,itemsep=0.45em]
\item \textbf{Exact problem predicate and public ordered form.}
  \path{RequestProject/Erdos287/ProblemStatement.lean}; declarations
  \path{Erdos287Counterexample} and
  \path{erdos287_seq_of_no_counterexample}.
\item \textbf{Finite blocker and verified range.}
  \path{Gap2CE.holes_isolated}, \path{Gap2CE.blocker_window}, and
  \path{no_Erdos287Counterexample_of_max_le_4e9} in the finite-range files.
\item \textbf{Large-$M$ conditional compiler.}
  \path{WindowPairSupply} and
  \path{no_Erdos287Counterexample_of_closure} in
  \path{RequestProject/Erdos287/ClosureInputs.lean}.
\item \textbf{Public analytic-source status through V14.}
  \path{Erdos287VaughanV14Status.lean} in the public status directory; exact structural
  spine banked, analytic child explicitly open.
\end{itemize}

\subsection*{Midpoint V15 formalization source -- public replay pending}
\begin{itemize}[leftmargin=1.6em,itemsep=0.45em]
\item \path{AffineMuLogIdentity.lean}: exact $\Lambda=\mu*\log$ identity and affine
  specialization.
\item \path{AffineMuLogHardSource.lean} and \path{AffineMuLogLine.lean}: exact
  small-$q$/small-$r$/hard partition and determinant-one line.
\item \path{BalancedSevenFinite.lean} and \path{BalancedSevenPolarization.lean}:
  coefficient $-20$, finite collision/exponent routers, and labelled polarization.
\item \path{PolarizedOmega7SignedEoD} and \path{MuLogComparisonLowCondMatch}:
  deliberately uninhabited open specifications, not analytic proofs.
\end{itemize}

\begin{thebibliography}{9}
\bibitem{ErdosGraham1980}
P. Erd\H{o}s and R. L. Graham,
\emph{Old and New Problems and Results in Combinatorial Number Theory},
Monographies de L'Enseignement Math\'ematique, 1980, p.~33.

\bibitem{ErdosProblems287}
T. F. Bloom,
\emph{Erd\H{o}s Problem \#287}, Erd\H{o}s Problems archive,
problem provenance and status page, accessed August 2026.

\bibitem{FordMaynard2024}
K. Ford and J. Maynard,
\emph{On the theory of prime producing sieves},
arXiv:2407.14368, 2024.

\bibitem{DGS2017}
S. Drappeau, A. Granville and X. Shao,
\emph{Smooth-supported multiplicative functions in arithmetic progressions beyond the
$x^{1/2}$-barrier}, Mathematika 63 (2017), 895--918.

\bibitem{Pascadi2025}
A. Pascadi,
\emph{On the exponents of distribution of primes and smooth numbers},
arXiv:2505.00653, 2025.

\bibitem{LeanMathlib}
The Lean and Mathlib communities,
\emph{Lean 4 and Mathlib}, formal proof assistant and mathematical library.
\end{thebibliography}

\end{document}
